\documentclass[10pt,a4paper]{amsart}
\usepackage[margin=19mm]{geometry}
\usepackage{amsmath,amssymb,mathtools}
\usepackage[scr=rsfs,cal=euler]{mathalfa}
\usepackage{xfrac}
\usepackage[unicode,hidelinks,bookmarks=false]{hyperref}
\newcommand{\ie}{i.e.\ }
\newcommand{\cf}{cf.\ }
\newcommand{\resp}{respectively\ }
\newcommand{\Subsection}{\subsection}
\providecommand{\nobreakdash}{\nobreak\hbox{-}}
\setlength{\emergencystretch}{2em}
\multlinegap0pt
\usepackage{amssymb}
\usepackage{dsfont}
\usepackage[shortlabels]{enumitem}
\usepackage{cancel}
\usepackage{tikz}
\newtheorem{thm}{Theorem}
\newtheorem{rem}{Remark}
\newcommand{\N}{\mathbb{N}}
\newcommand{\abs}[1]{\left\vert#1\right\vert}
\newcommand{\cl}{\operatorname{cl}}
\newcommand{\dimf}{\operatorname{dim}_{f}}
\newcommand{\dist}{\operatorname{dist}}
\newcommand{\eps}{\varepsilon}
\renewcommand{\geq}{\geqslant}
\renewcommand{\leq}{\leqslant}
\title{A panoramic view of exponential attractors}
\author{Radoslaw Czaja, Stefanie Sonner}
\date{Source: arXiv:2506.10388v2 (CC BY 4.0)}
\begin{document}
\maketitle

\noindent\emph{Source and licence.} A panoramic view of exponential attractors. Version arXiv:2506.10388v2 (CC BY 4.0), distributed by arXiv under the Creative Commons Attribution 4.0 International licence, http://creativecommons.org/licenses/by/4.0/, as recorded in the arXiv licence metadata for this exact version and shown on the abstract page https://arxiv.org/abs/2506.10388v2. Full licence notice: \texttt{licenses/arxiv-2506.10388-CC-BY-4.0.txt}.\par\smallskip

\noindent\textit{Editorial scope.} Counted unit: the article's thm thm:SEM3 (source0.tex lines 1539--1563). The article's complete original proof of it is delivered with it (source0.tex lines 1565--1721, one \texttt{proof} environment of 157 lines). No mathematics is written, repaired or re-derived: the statement and the proof are the article's own lines, in the article's own notation, and no retained line is substituted or re-flowed.  Retained uncounted as the source-local context the proof needs: lines 343--364; lines 351--353; lines 367--375.  Nothing was omitted from any retained span, so the package carries no omission record.  No retained line contains a citation, so the wrapper carries no bibliography and no bibliography entry.  No retained line contains a figure environment, an image-inclusion command or drawing code, so no image is delivered and the package declares no asset dependency.  Editorial formatting only, in addition: an AMS \texttt{amsart} preamble that restates the article's own declarations and macros used by the retained lines, namely the environments \texttt{thm}, \texttt{rem} on separate counters, together with the standard packages those lines need.  The retained environments print in this document as Theorem 1, Remark 1 in order of appearance, whereas the article prints the same items with the numbering of its own full document.  The complete native source of the article as distributed in the arXiv e-print for this exact version is bundled unchanged as \texttt{sources/179726/source0.tex}.
\begin{thm}\label{thm:EXDEXPAT}
Let $\{S(t)\colon t\geq 0\}$ be an asymptotically closed semigroup on a complete metric space $(V,d)$ and let $T>0$.
Then, the following statements are equivalent:
\begin{itemize}
\item[$(1)$] There exists a $T$-discrete exponential attractor $\mathbf{M_0}$ in $V$ for the semigroup.
\item[$(2)$] There exists a nonempty bounded 
absorbing set $\mathbf{B}\subseteq V$ for the semigroup
such that the covering condition 
\begin{equation}\label{e:CONDBALLSSEMB}
N^{V}(S(kT)\mathbf{B},aq^{k})\leq b h^{k},\ k\in\N,\ k\geq k_0,
\end{equation}
holds for some $k_0\in\N$, $a,b>0$, $q\in(0,1)$ and $h\geq 1$. 
\end{itemize}
Moreover, if the covering condition \eqref{e:CONDBALLSSEMB} holds, then
$$\mathbf{M_0}=\mathbf{A}\cup\mathbf{E_0}=\cl_{V}\mathbf{E_0}\subseteq \mathbf{B},$$
is a $T$-discrete exponential attractor with rate of attraction
$\xi\in(0,\frac{1}{T}\ln{\frac{1}{q}})$, and its fractal dimension is bounded by
\begin{equation}\label{e:ESTDIMM3}
\dimf^{V}(\mathbf{M_0})\leq\log_{\frac{1}{q}}h.
\end{equation}
Here, $\mathbf{E_0}$ is a certain countable subset of $\mathbf{B}$ and $\mathbf{A}=\Lambda^{V}(\cl_{V}\mathbf{E_0})$ is the global attractor for the semigroup.
\end{thm}

\begin{equation}
N^{V}(S(kT)\mathbf{B},aq^{k})\leq b h^{k},\ k\in\N,\ k\geq k_0,
\end{equation}

\begin{rem}\label{rem:POSINV}
Without loss of generality, in Theorem~\ref{thm:EXDEXPAT} we can assume that the absorbing set $\mathbf{B}$ is \emph{positively invariant}.
Indeed, if $B_0$ is a~bounded absorbing set that satisfies the covering condition \eqref{e:CONDBALLSSEMB} 
and $S(t)B_0\subseteq B_0$ for all $t\geq t_{B_0}$, $t_{B_0}\geq0$ being the absorbing time for $B_0$, then the positively invariant subset
\begin{equation*}
\mathbf{B}=\bigcup_{t\geq t_{B_0}}S(t)B_0\subseteq B_0
\end{equation*}
is a bounded absorbing set that satisfies \eqref{e:CONDBALLSSEMB} with a possibly  larger constant $a>0$.
\end{rem}

\begin{thm}\label{thm:SEM3}
Let $\{S(t)\colon t\geq 0\}$ be an asymptotically closed semigroup on a~complete metric space $(V,d)$, $\mathbf{B}$ be a bounded absorbing set, $T>0$ and assume
that the covering condition \eqref{e:CONDBALLSSEMB}
holds with some $k_0\in\N$, $q\in(0,1)$, $a,b>0$ and $h\geq 1$.
Assume further that there exist $T_2>T_1\geq 0$,
$\zeta>0$ and $\nu>0$ such that
\begin{equation}\label{e:TIMEHOLDER}
d(S(t_1)x,S(t_2)x)\leq\zeta\abs{t_1-t_2}^{\nu},\ t_1,t_2\in[T_1,T_2],\ x\in\mathbf{B}.
\end{equation}
If $T_1>0$ then we also assume that for some $N\in\N$ such that $NT\geq T_1$
\begin{equation}\label{e:LIPTSEMN}
d(S(NT)x,S(NT)y)\leq L_{N}d(x,y),\ x,y\in \mathbf{B},
\end{equation}
holds with some $L_{N}\geq 0$. Then there exists an exponential attractor
$\mathbf{M}$ (independent of $\nu$, $\zeta$, $T_2$) with rate of attraction $\xi\in(0,\frac{1}{T}\ln\frac{1}{q})$
and its fractal dimension is bounded by
\begin{equation*}
\dimf^{V}(\mathbf{M})\leq\tfrac{1}{\nu}+\log_{\frac{1}{q}}h.
\end{equation*}
Moreover,
\begin{equation*}
\mathbf{M}=\mathbf{A}\cup\mathbf{E}=\cl_{V}\mathbf{E}\subseteq\mathbf{B},
\end{equation*}
where $\mathbf{E}=\bigcup_{t\in I}S(t)\mathbf{E_0} \subseteq\mathbf{B}$ for some compact interval $I$, $\mathbf{E_0}$ is the set constructed in Theorem~\ref{thm:EXDEXPAT} and $\mathbf{A}=\Lambda^{V}(\cl_{V}\mathbf{E_0})$ is the global attractor for the semigroup.
\end{thm}

\begin{proof}
\textbf{Step 1.} Without loss of generality we can assume that $\mathbf{B}$ is positively invariant by Remark~\ref{rem:POSINV}. We first show that there exists a~set $\mathbf{E}$ (independent of $\nu$, $\zeta$, $T_2$) with the following properties:
\begin{itemize}
\item[(i)] $\mathbf{E}\subseteq\mathbf{B}$, $\mathbf{E}$ is precompact in $V$,
\item[(ii)] $S(t)\mathbf{E}\subseteq\mathbf{E},\ t\geq 0,$
\item[(iii)] the fractal dimension of $\mathbf{E}$ is bounded by
\begin{equation*}
\dimf^{V}(\mathbf{E})\leq\tfrac{1}{\nu}+\log_{\frac{1}{q}}h,
\end{equation*}
\item[(iv)] for any $\xi\in(0,\frac{1}{T}\ln\frac{1}{q})$,
and any bounded subset $G$ of $V$ we have
\begin{equation}\label{e:EXPATESEM2}
\lim_{t\rightarrow\infty}e^{\xi t}\dist^{V}(S(t)G,\mathbf{E})=0.
\end{equation}
\end{itemize}
Indeed, if $T_1=0$ then we set $N=0$, otherwise let $N\in\N$ be such that $NT\geq T_1$
and \eqref{e:LIPTSEMN} holds. We define
$$\mathbf{E}=\bigcup_{p\in[NT,(N+1)T]}S(p)\mathbf{E_0},$$
where $\mathbf{E_0}$, constructed in Step 2 in the proof of Theorem~\ref{thm:EXDEXPAT}, satisfies ($e_1$), ($e_2$) from there. Since $\mathbf{E_0}\subseteq \mathbf{B}$
and $\mathbf{B}$ is positively invariant, we get $\mathbf{E}\subseteq \mathbf{B}$. For $t\geq 0$ we observe that
\begin{align*}
S(t)\mathbf{E}&=\bigcup_{p\in[NT,(N+1)T]}S(p+t)\mathbf{E_0}\subseteq\bigcup_{s\in[0,T)}\bigcup_{l\in\N_0}S((N+l)T+s)\mathbf{E_0}\\
&\subseteq\bigcup_{s\in[0,T)}S(NT+s)\mathbf{E_0}\subseteq\mathbf{E},
\end{align*}
which shows (ii).
Let $\xi\in(0,\frac{1}{T}\ln\frac{1}{q})$ and $G\subseteq V$ be bounded.
If $T_1=0$ then $\mathbf{E_0}\subseteq\mathbf{E}$ and \eqref{e:EXPATESEM2} follows directly from
($e_2$). Otherwise, \eqref{e:LIPTSEMN} is assumed and we
know that $S(t)G\subseteq\mathbf{B}$ for $t\geq t_{G}$ and for any $\eps>0$ there exists $t_\eps\geq 0$
such that
$$e^{\xi t}\dist^{V}(S(t)G,\mathbf{E_0})<\frac{\eps}{e^{\xi NT}(L_{N}+1)},\ t\geq t_\eps.$$
Fix $\eps>0$ and let $t\geq NT+t_\eps+t_{G}$. Then we have by \eqref{e:LIPTSEMN}
\begin{align*}
e^{\xi t}\dist^{V}(S(t)G,\mathbf{E})&\leq e^{\xi t}\dist^{V}(S(NT)S(t-NT)G,S(NT)\mathbf{E_0})\\
&\leq e^{\xi NT}L_{N}e^{\xi(t-NT)}\dist^{V}(S(t-NT)G,\mathbf{E_0})<\eps.
\end{align*}
We are left to prove (iii). Now let $\tau=T_2-T_1>0$. First we show that
\begin{equation}\label{e:HOLDERONSMALL}
d(S(t_1)x,S(t_2)x)\leq\zeta\abs{t_1-t_2}^{\nu},\ x\in \mathbf{B},\ t_1,t_2\in[T_1+l\tau,T_1+(l+1)\tau],\ l\in\N_0.
\end{equation}
Indeed, let $l\in\N_0$, $t_1,t_2\in[T_1+l\tau,T_1+(l+1)\tau]$ and $x\in\mathbf{B}$. Then $t_i=T_1+l\tau+s_i$, $s_i\in[0,\tau]$, $i=1,2$, and
$$d(S(t_1)x,S(t_2)x)=d(S(T_1+s_1)S(l\tau)x,S(T_1+s_2)S(l\tau)x)\leq\zeta\abs{t_1-t_2}^{\nu}.$$
By ($e_1$) from Step 2 in the proof of Theorem~\ref{thm:EXDEXPAT} we also have
$$\mathbf{E}=\bigcup_{k\geq k_0}\bigcup_{p\in[NT,(N+1)T]}S(p)Q_{k},$$
where $Q_k\subseteq S(kT)\mathbf{B}$ and $\displaystyle\# Q_k\leq b\sum_{j=0}^{k-k_0}h^{k-j}$.

Let $m_0\geq k_0$ be such that
\begin{equation}\label{e:M0}
aq^{m}\leq 2\zeta\tau^{\nu}\text{ for }m\geq m_0.
\end{equation}
Let $m\geq m_0$ and $p\in[NT,(N+1)T]$. For $k\geq m$ we have
$$S(p)Q_k\subseteq S(p)S(kT)\mathbf{B}=S(mT)S((k-m)T+p)\mathbf{B}\subseteq S(mT)\mathbf{B},$$
and thus
$$\mathbf{E}\subseteq\bigcup_{k=k_0}^{m}\bigcup_{p\in[NT,(N+1)T]}S(p)Q_k\cup S(mT)\mathbf{B},\ m\geq m_0.$$
We now construct a finite covering of this set
by balls with centers in $\mathbf{E}$ and radii $2aq^{m}$. 
To this end, we denote the elements of $\displaystyle\bigcup_{k=k_0}^{m}Q_k$ by $\{x_i\colon i=1,\ldots,i_m\}$ and observe that 
$$i_{m}=\# \bigcup_{k=k_0}^{m}Q_k\leq\sum_{k=k_0}^{m}\# Q_k\leq b\sum_{k=k_0}^{m}\sum_{j=0}^{k-k_0}h^{k-j}\leq b(m-k_0+1)^{2}h^{m}.$$
The interval $[NT,(N+1)T]$ can be covered by $\left[\frac{T}{\tau}\right]+2$ intervals of length $\tau$
of the form $[T_1+l\tau,T_1+(l+1)\tau]$ with $l\in\N_0$. 
We further subdivide each such interval into intervals of length
$\left(\frac{aq^{m}}{2\zeta}\right)^{\frac{1}{\nu}}\leq \tau$ 
by \eqref{e:M0}, plus possibly one interval of smaller length. We note that 
$\Big[\tau\Big(\frac{2\zeta}{aq^{m}}\Big)^{\frac{1}{\nu}}\Big]+1$ such intervals cover 
$[T_1+l\tau,T_1+(l+1)\tau]$, since
$\Big(\frac{aq^{m}}{2\zeta}\Big)^{\frac{1}{\nu}}\Big(\Big[\tau\Big(\frac{2\zeta}{aq^{m}}\Big)^{\frac{1}{\nu}}\Big]+1\Big)>\tau$.
Let $I_j$, $j=1,\ldots, j_m$, be an arbitrary one of these intervals, where
$$j_m\leq\left(\left[\tfrac{T}{\tau}\right]+2\right)\Big(\Big[\tau\Big(\tfrac{2\zeta}{aq^{m}}\Big)^{\frac{1}{\nu}}\Big]+1\Big).$$
We choose $p_j\in I_j$. Then, for any $p\in I_j$ we get by \eqref{e:HOLDERONSMALL}
$$d(S(p)x_i,S(p_j)x_i)\leq\zeta\abs{p-p_j}^{\nu}\leq\frac{1}{2}aq^{m}<aq^{m},$$
and thus, 
$$\bigcup_{i=1}^{i_m}\bigcup_{p\in[NT,(N+1)T]}S(p)x_i\subseteq\bigcup_{i=1}^{i_m}\bigcup_{j=1}^{j_m}B^{V}(S(p_j)x_i,aq^{m}).$$
This implies that
$$\bigcup_{k=k_0}^{m}\bigcup_{p\in[NT,(N+1)T]}S(p)Q_k\cap\mathbf{E}\subseteq\bigcup_{i=1}^{i_m}\bigcup_{j=1}^{j_m}B^{V}(u_{i,j},2aq^{m})$$
for some $u_{i,j}\in\mathbf{E}$. Using \eqref{e:CONDBALLSSEMB} and \eqref{e:M0} we obtain for $m\geq m_0$
\begin{equation*}
\begin{split}
N^{V}(\mathbf{E},2aq^{m})\leq i_m j_m+bh^{m}&\leq b(m-k_0+1)^{2}h^{m}\left(\left[\tfrac{T}{\tau}\right]+2\right)\left(\tau\left(\tfrac{2\zeta}{aq^{m}}\right)^{\frac{1}{\nu}}+1\right)+bh^{m}\\
&\leq c(m-k_0+1)^{2}h^{m}\left(\tfrac{2\zeta}{aq^{m}}\right)^{\frac{1}{\nu}},
\end{split}
\end{equation*}
where $c$ is a positive constant independent of $m$.
Since $2aq^m$ converges to $0$ as $m\to\infty$, it follows that $\mathbf{E}$
is precompact in $V$. Consider an arbitrary sequence $\eps_n>0$ converging to $0$ and choose
$m_n\in\N$ such that $m_n\geq m_0$ and
$$2aq^{m_n}\leq \eps_n<2aq^{m_n-1}<1\text{ for }n\geq n_0.$$
Then, we have
$$\log_{\frac{1}{\eps_n}}N^{V}(\mathbf{E},\eps_n)\leq\frac{\ln{c}+2\ln(m_n-k_0+1)+m_n\ln{h}+\frac{1}{\nu}\left(\ln(\frac{2\zeta}{a})-m_n\ln{q}\right)}{-\ln(2a)-(m_n-1)\ln{q}},$$
which implies that
$$\dimf^{V}(\mathbf{E})\leq\frac{\ln{h}-\tfrac{1}{\nu}\ln{q}}{-\ln{q}}=\tfrac{1}{\nu}+\log_{\frac{1}{q}}h$$
and completes the proof of the claims (i)--(iv).

\textbf{Step 2.}
By Theorem~\ref{thm:EXDEXPAT} the global attractor for the semigroup exists, $\mathbf{A}=\Lambda^{V}(\cl_{V}\mathbf{E_0})$  and
\begin{equation*}
\dimf^{V}(\mathbf{A})\leq\log_{\frac{1}{q}}h.
\end{equation*}
We define
$$\mathbf{M}=\mathbf{A}\cup\mathbf{E}\subseteq\mathbf{B}$$
and show that it is an exponential attractor. 

The set $\mathbf{M}$ is nonempty, positively invariant, precompact and
$$\dimf^{V}(\mathbf{M})=\max\{\dimf^{V}(\mathbf{A}),\dimf^{V}(\mathbf{E})\}\leq\tfrac{1}{\nu}+\log_{\frac{1}{q}}h.$$
Moreover,  for any $\xi\in(0,\frac{1}{T}\ln{\frac{1}{q}})$, and
any bounded subset $G$ of $V$ we have
$$\lim_{t\to\infty}e^{\xi t}\dist^{V}(S(t)G,\mathbf{M})=0.$$
Hence, it remains to show that $\mathbf{M}$ is compact.

Consider a sequence $x_n\in\mathbf{M}$, $n\in\N$. If infinitely many of its elements belong
to $\mathbf{A}$, then by the compactness of $\mathbf{A}$, there exists a subsequence convergent
to an element of $\mathbf{A}\subseteq\mathbf{M}$. Otherwise,  there is a subsequence
$$x_{n_j}\in\mathbf{E}=\bigcup_{k\geq k_0}\bigcup_{p\in[NT,(N+1)T]}S(p)Q_k.$$
Thus $x_{n_j}=S(p_j)y_j$, where $p_{j}\in[NT,(N+1)T]$, $y_j\in Q_{k_j}$ for some $k_j\geq k_0$.
Taking a~subsequence if necessary, we can assume that $p_j\to p_0\in[NT,(N+1)T]\subset[T_1,\infty)$
by the compactness of the interval.

If $p_0\in(T_1+l\tau,T_1+(l+1)\tau)$ for some $l\in\N_0$, then $p_j\in(T_1+l\tau,T_1+(l+1)\tau)$
for large $j$. Moreover, by \eqref{e:HOLDERONSMALL} we have
$$d(S(p_j)x,S(p_0)x)\leq\zeta\abs{p_j-p_0}^{\nu},\ x\in\mathbf{B}.$$
If $p_0=T_1+l\tau$ for some $l\in\N$, then $p_j\in(T_1+l\tau-\frac{\tau}{2},T_1+l\tau+\frac{\tau}{2})$
for large $j$. Moreover, for $x\in\mathbf{B}$ and $t_1,t_2\in(T_1+l\tau-\frac{\tau}{2},T_1+l\tau+\frac{\tau}{2})$
we have $t_i=T_1+l\tau-\frac{\tau}{2}+s_i$, $s_i\in(0,\tau)$, $i=1,2$ and by \eqref{e:TIMEHOLDER}
and the positive invariance of $\mathbf{B}$ we get
\begin{equation*}
\begin{split}
d(S(t_1)x,S(t_2)x)&=d(S(T_1+s_1)S(l\tau-\tfrac{\tau}{2})x,S(T_1+s_2)S(lT-\tfrac{\tau}{2})x)\\
&\leq\zeta\abs{s_1-s_2}^{\nu}=\zeta\abs{t_1-t_2}^{\nu}.
\end{split}
\end{equation*}
If $p_0=T_1=NT$, then we directly apply \eqref{e:TIMEHOLDER} and conclude from
the above considerations that
\begin{equation*}
S(p_j)x\to S(p_0)x\text{ as }j\to\infty\text{ for any }x\in\mathbf{B}.
\end{equation*}
We distinguish two cases.
If $K=\sup\{k_j\colon j\in\N\}<\infty$, then the elements $y_j$ belong to the finite set $\bigcup_{k=k_0}^{K}Q_k$
and we find a constant subsequence $y_{j_l}=y_0\in\bigcup_{k=k_0}^{K}Q_k\subseteq\mathbf{B}$.
It follows that 
$$x_{n_{j_l}}=S(p_{j_l})y_{j_l}=S(p_{j_l})y_0\to S(p_0)y_0\in\bigcup_{k=k_0}^{K}\bigcup_{p\in[NT,(N+1)T]}S(p)Q_k\subseteq\mathbf{E}\subseteq\mathbf{M}.$$
If $\sup\{k_j\colon j\in\N\}=\infty$, then there exists a~subsequence $k_{j_l}\to\infty$ such that
$$x_{n_{j_l}}=S(p_{j_l})y_{j_l}\in S(p_{j_l})S(k_{j_l}T)\mathbf{B}\subseteq S(k_{j_l}T)\mathbf{B}.$$
Since $\dist^{V}(S(k_{j_l}T)\mathbf{B},\mathbf{A})\to 0$ as $l\to\infty$, the sequence $x_{n_{j_l}}$
has a convergent subsequence to an element of $\mathbf{A}\subseteq\mathbf{M}$.

We now show that $\mathbf{A}\cup\mathbf{E}=\cl_{V}\mathbf{E}$.
Since $\cl_{V}\mathbf{E}$ is a compact set attracting all bounded sets, by
the minimality of the global attractor we have $\mathbf{A}\subseteq\cl_{V}\mathbf{E}$ and thus $\mathbf{A}\cup\mathbf{E}\subseteq\cl_{V}\mathbf{E}$.
For the converse inclusion, let $x\in\cl_{V}\mathbf{E}$. Then there exists a sequence
$x_j\in\mathbf{E}$ such that $x_j\to x$. Therefore, there are $k_j\geq k_0$
and $p_j\in[NT,(N+1)T]$ such that $x_j=S(p_j)y_j$ with some $y_j\in Q_{k_j}$.
Taking a subsequence if necessary, we assume that $p_j\to p_0\in[NT,(N+1)T]$.
Following the above arguments, if $K=\sup\{k_j\colon j\in\N\}<\infty$, then
$x_j$ has a convergent subsequence to an element of $\mathbf{E}$ and hence,  $x\in\mathbf{E}$.
If $\sup\{k_j\colon j\in\N\}=\infty$, then $x_j$ has a~convergent subsequence to an element
of $\mathbf{A}$, so $x\in\mathbf{A}$. We conclude that $x\in\mathbf{A}\cup\mathbf{E}$,
which completes the proof.
\end{proof}

\end{document}
